\chapter{Mesures du risque de marché}\label{ch:rc:market-risk-measures}

En octobre 1994, J.P. Morgan publie RiskMetrics~: une méthode pour transformer les positions d'une société de trading
en un seul nombre, la perte qu'elle ne doit pas dépasser sur plus d'un jour sur vingt (ou
sur cent), accompagnée des volatilités quotidiennes et des corrélations nécessaires à son calcul,
le tout gratuitement. Deux ans plus tard, le Comité de Bâle autorise les banques à fixer leurs fonds propres
au titre du risque de marché à partir de leur propre version de ce nombre, à condition
de le confronter chaque jour à ce qui s'est réellement produit. La \omterm{def:rc:market-risk-measures:var}{valeur en risque} devient le langage
des limites de risque, des fonds propres et de la communication financière. Ses faiblesses deviennent tout aussi connues~: elle ne dit rien de l'ampleur des pertes au-delà du seuil,
elle peut récompenser la concentration, et ses estimations sont en retard sur les marchés qu'elle mesure. Ce chapitre
ouvre la partie~IV en calculant ce nombre de trois façons sur un portefeuille réel et des données réelles,
en le remplaçant par l'\omterm{def:rc:market-risk-measures:es}{expected shortfall}, en soumettant les deux à un backtest, puis en les allouant à leurs sources.

\section{Distributions des pertes et valeur en risque}

\begin{definition}[Facteur de risque]\label{def:rc:market-risk-measures:factor}
Un \emph{facteur de risque}\index{facteur de risque} est une variable de marché dont les variations font
varier la valeur des positions~: un rendement à une échéance donnée, un taux de change, une volatilité implicite, un spread de crédit.
Un modèle de risque décrit la distribution conjointe des variations des facteurs de risque sur un horizon
et réévalue les positions sous ces variations.
\end{definition}

\begin{definition}[Valeur en risque]\label{def:rc:market-risk-measures:var}
La \emph{valeur en risque}\index{valeur en risque} d'un portefeuille au niveau de confiance $\alpha$ sur
un horizon $h$ est le quantile d'ordre $\alpha$ de sa perte $L = -\Delta V$ sur $h$~:
$\mathrm{VaR}_\alpha = \inf\{\ell:\P(L>\ell)\le1-\alpha\}$. À 99~\% sur un jour, c'est une perte dépassée
un jour sur cent.
\end{definition}

\begin{example}[Le portefeuille]\label{ex:rc:market-risk-measures:book}
Un petit portefeuille en USD au 23 septembre 2026~: une position longue de USD~200 millions d'une obligation du Trésor à dix ans au pair,
une position courte de USD~300 millions d'une obligation à deux ans, une position longue de EUR~100 millions, une position courte de JPY~5 milliards, et une position courte
d'un straddle EURUSD à trois mois, à la monnaie, sur EUR~400 millions (volatilité implicite de 8~\%, maintenue
fixe). Ses quatre \omterm{def:rc:market-risk-measures:factor}{facteurs de risque} sont les variations quotidiennes des rendements au pair à deux et dix ans (4.85~\%
et 5.11~\% ce jour-là) et les rendements logarithmiques de l'EURUSD (1.1411) et de l'USDJPY (157.92). Leurs volatilités
quotidiennes EWMA sont de 6.4 et 5.7 points de base, soit 0.28~\% et 0.65~\%.
\end{example}

\section{La valeur en risque de trois façons}

\begin{definition}[Simulation historique]\label{def:rc:market-risk-measures:hs}
La \emph{simulation historique}\index{simulation historique} applique chacune des $n$ dernières
variations observées des \omterm{def:rc:market-risk-measures:factor}{facteurs de risque} aux positions d'aujourd'hui, les réévalue et lit la VaR sur
la distribution empirique des $n$ P\&L~: avec $k = \lceil(1-\alpha)n\rceil$, la $k$-ième plus grande perte.
\end{definition}

\begin{definition}[VaR paramétrique]\label{def:rc:market-risk-measures:param}
La \emph{VaR paramétrique}\index{VaR paramétrique} (delta-normale) linéarise
le P\&L en fonction des variations des facteurs, $\Delta V\approx\delta^\top x$, suppose $x$ normal de covariance $\Sigma$,
et donne $\mathrm{VaR}_\alpha = z_\alpha\sqrt{\delta^\top\Sigma\delta}$. Avec des covariances à pondération exponentielle
($\Sigma_t = \lambda\Sigma_{t-1}+(1-\lambda)x_tx_t^\top$, $\lambda = 0.94$ pour les données quotidiennes de RiskMetrics), elle
suit rapidement la volatilité.
\end{definition}

\begin{definition}[VaR Monte Carlo]\label{def:rc:market-risk-measures:mc}
La \emph{VaR Monte Carlo}\index{VaR Monte Carlo} simule les variations des facteurs à partir d'un modèle
(ici normal, de covariance EWMA), réévalue le portefeuille en entier à chaque tirage
et lit le quantile.
\end{definition}

\begin{definition}[Simulation historique filtrée]\label{def:rc:market-risk-measures:fhs}
La \emph{simulation historique filtrée}\index{simulation historique filtrée} remet à l'échelle chaque variation passée
d'un facteur par le rapport entre l'estimation de volatilité d'aujourd'hui et l'estimation de son
jour~: $\tilde x_s = x_s\,\hat\sigma_{\text{aujourd'hui}}/\hat\sigma_s$, avant de l'appliquer~: elle conserve les formes historiques
et la dépendance des mouvements, ainsi que le niveau de volatilité d'aujourd'hui.
\end{definition}

\begin{definition}[Approximation delta--gamma]\label{def:rc:market-risk-measures:dg}
L'\emph{approximation delta--gamma}\index{approximation delta--gamma} remplace la réévaluation complète par
le développement du second ordre $\Delta V\approx\delta^\top x+\tfrac12x^\top\Gamma x$, avec les dérivées
première et seconde du portefeuille par rapport aux \omterm{def:rc:market-risk-measures:factor}{facteurs de risque}.
\end{definition}

\begin{example}[Un portefeuille, plusieurs chiffres]\label{ex:rc:market-risk-measures:methods}
Sur les 500 derniers jours, la \omterm{def:rc:market-risk-measures:hs}{simulation historique} donne une VaR à un jour et à 99~\% de USD~2.10 millions~;
la VaR delta-normale EWMA vaut 1.58 million, celle à pondération égale 1.60 million, la méthode Monte Carlo
avec réévaluation complète sur la covariance EWMA 1.66 million, et la \omterm{def:rc:market-risk-measures:fhs}{simulation historique filtrée} 2.09 millions
(\cref{fig:rc:market-risk-measures:methods}). La \omterm{def:rc:market-risk-measures:hs}{simulation historique} avec réévaluation delta--gamma donne 2.10 millions,
le delta seul 1.91~: le gamma du straddle court ajoute des pertes dans les deux queues.
Les méthodes normales donnent des chiffres plus faibles parce que les mouvements historiques ont des queues plus épaisses
que celles de la loi normale.
\end{example}

\omcode{code/firm/varmodel/firm_varmodel.py}{23}{50}{VaR et ES historiques, covariance et volatilité EWMA, et scénarios historiques filtrés.}\label{lst:rc:market-risk-measures:var}

\begin{definition}[Règle de la racine carrée du temps]\label{def:rc:market-risk-measures:sqrt}
La \emph{règle de la racine carrée du temps}\index{règle de la racine carrée du temps} ramène une VaR à un jour à $h$
jours en la multipliant par $\sqrt h$~: exact pour des variations normales indépendantes et identiquement distribuées, de moyenne nulle,
avec des positions fixes, et une approximation dans les autres cas.
\end{definition}

\section{Expected shortfall et cohérence}

\begin{definition}[Expected shortfall]\label{def:rc:market-risk-measures:es}
L'\emph{expected shortfall}\index{expected shortfall} de niveau $\alpha$ est la perte moyenne dans les pires $1-\alpha$ des
issues, $\mathrm{ES}_\alpha = \frac1{1-\alpha}\int_\alpha^1\mathrm{VaR}_u\,du$~; en historique, c'est la moyenne
des $k$ plus grandes pertes.
\end{definition}

\begin{definition}[Mesure de risque cohérente]\label{def:rc:market-risk-measures:coherent}
Une mesure de risque $\rho$ est une \emph{mesure de risque cohérente}\index{mesure de risque cohérente} si elle est
monotone, invariante par translation ($\rho(L+c) = \rho(L)+c$), positivement homogène ($\rho(\lambda L) = \lambda\rho(L)$
pour $\lambda\ge0$) et sous-additive ($\rho(L_1+L_2)\le\rho(L_1)+\rho(L_2)$). L'\omterm{def:rc:market-risk-measures:es}{expected shortfall} est cohérente~;
la VaR n'est pas sous-additive en général.
\end{definition}

\begin{example}[Deux obligations concentrées]\label{ex:rc:market-risk-measures:subadd}
Chacune de deux obligations fait défaut indépendamment avec une probabilité de 0.9~\%, en perdant 100. Prise isolément, chacune
a une VaR à 99~\% nulle~; ensemble, la probabilité d'au moins un défaut vaut 1.79~\%, donc la VaR
à 99~\% vaut 100~: la diversification a fait monter la VaR. L'\omterm{def:rc:market-risk-measures:es}{expected shortfall} à 99~\% vaut 90 pour chaque obligation et 100.8 pour
la paire, soit moins que leur somme de 180.
\end{example}

Pour un P\&L normal, l'\omterm{def:rc:market-risk-measures:es}{expected shortfall} à 97.5~\% égale presque exactement la VaR à 99~\% (2.338 contre 2.326 écarts
types), ce qui explique pourquoi le Comité de Bâle a retenu 97.5~\% lorsqu'il a remplacé la VaR par l'ES dans sa
norme 2019 sur le risque de marché. Sur le livre, l'ES EWMA à 97.5~\% vaut 1.59~million de dollars US, à côté de la
VaR de 1.58~; l'ES historique vaut 2.34~million de dollars US contre une VaR historique de 2.10~: les queues épaisses apparaissent
d'abord dans l'ES.

\begin{omfigure}
\begin{tikzpicture}
\begin{axis}[width=11.5cm,height=5.4cm,ybar,bar width=9pt,
  symbolic x coords={HS,EWMA,equal-wt,MC,FHS,HS delta-gamma},xtick=data,xticklabels={HS,EWMA,{équi-\\pondéré},MC,FHS,{HS\\delta-gamma}},x tick label style={align=center},
  ylabel={millions de dollars},
  tick label style={font=\small},label style={font=\small},
  ymin=0,ymax=2.8,grid=major,grid style={gray!25},
  legend columns=2,legend style={font=\footnotesize,at={(0.5,-0.3)},anchor=north,draw=none,fill=none,/tikz/every even column/.append style={column sep=8pt}},
  table/col sep=comma]
\addplot[fill=omDef!60,draw=omDef] table[x=method,y=var99]{figdata/rates-credit-risk/21-market-risk-measures/methods.csv};
\addplot[fill=omThm!60,draw=omThm] table[x=method,y=es975]{figdata/rates-credit-risk/21-market-risk-measures/methods.csv};
\legend{VaR à 99~\%,ES à 97.5~\%}
\end{axis}
\end{tikzpicture}

\omcaption{VaR à 99~\% et ES à 97.5~\% sur un jour du livre au 23 septembre 2026, par méthode. Les méthodes normales
concordent entre elles et donnent des valeurs plus basses~; la \omterm{def:rc:market-risk-measures:hs}{simulation historique} et la \omterm{def:rc:market-risk-measures:fhs}{simulation historique filtrée}
voient les queues épaisses, et leur ES dépasse leur VaR. Données~: taux au pair du Trésor américain,
taux de référence de la BCE~; tutoriel du chapitre.}\label{fig:rc:market-risk-measures:methods}
\end{omfigure}

\begin{omfigure}
\begin{tikzpicture}
\begin{axis}[width=11.5cm,height=5.2cm,ybar,bar width=4pt,
  xlabel={P\&L d'un jour sous les scénarios historiques (millions de dollars)},ylabel={jours},
  tick label style={font=\small},label style={font=\small},
  xmin=-5.5,xmax=2.5,ymin=0,ymax=75,grid=major,grid style={gray!25},
  table/col sep=comma]
\addplot[fill=omDef!40,draw=omDef] table[x=pnl,y=count]{figdata/rates-credit-risk/21-market-risk-measures/hist.csv};
\draw[omThm,very thick] (axis cs:-2.097,0) -- (axis cs:-2.097,40) node[above,font=\footnotesize] {VaR};
\draw[omThm,very thick,dashed] (axis cs:-2.335,0) -- (axis cs:-2.335,28) node[above left,font=\footnotesize] {ES};
\end{axis}
\end{tikzpicture}

\omcaption{Distribution du P\&L du livre sur les 500 derniers jours historiques appliqués aux
positions actuelles, avec la VaR à 99~\% et l'ES à 97.5~\%. Données~: Trésor américain, BCE~; le
tutoriel du chapitre.}\label{fig:rc:market-risk-measures:hist}
\end{omfigure}

\section{Backtester la mesure}

\begin{definition}[Dépassement de VaR]\label{def:rc:market-risk-measures:exception}
Un \emph{dépassement de VaR}\index{dépassement de VaR} est un jour où la perte réalisée (ou hypothétique)
dépasse la VaR calculée la veille. Un modèle correct à 99~\% présente des dépassements sur 1~\% des jours,
de manière indépendante.
\end{definition}

\begin{definition}[Test de Kupiec, test de Christoffersen]\label{def:rc:market-risk-measures:tests}
Le \emph{test de Kupiec}\index{test de Kupiec} compare le nombre de dépassements $x$ sur $n$ jours à
le nombre attendu $pn$ par le rapport de vraisemblance $-2\ln[(1-p)^{n-x}p^x/((1-\hat p)^{n-x}\hat p^x)]$, $\hat p = x/n$,
suivant une loi du khi-deux à un degré de liberté. Le \emph{test de Christoffersen}\index{test de Christoffersen}
vérifie l'indépendance~: il compare la probabilité d'un dépassement après un dépassement avec la
probabilité après un jour calme.
\end{definition}

\begin{definition}[Test des feux tricolores]\label{def:rc:market-risk-measures:traffic}
Le \emph{test des feux tricolores}\index{test des feux tricolores} classe un modèle de VaR à 99~\% selon ses
dépassements sur 250 jours~: vert pour 0 à 4, jaune pour 5 à 9 avec des hausses de 0.40, 0.50,
0.65, 0.75 et 0.85 du facteur multiplicatif de capital de 3, rouge pour 10 ou plus (une hausse
de 1). Selon les règles de 1996, le capital est le plus élevé entre la VaR à dix jours de la veille et le facteur multiplié
par sa moyenne sur les soixante derniers jours ouvrables précédents.
\end{definition}

\omcode{code/firm/varmodel/firm_varmodel.py}{77}{108}{Tests du rapport de vraisemblance de Kupiec et de Christoffersen, et feux tricolores de Bâle.}\label{lst:rc:market-risk-measures:backtest}

\begin{example}[Quatre ans de backtest]\label{ex:rc:market-risk-measures:backtest}
On recalcule chaque modèle chaque jour du 2 septembre 2022 au 23 septembre 2026, sur le même livre
réévalué aux cours de chaque jour, face au P\&L du lendemain~: la \omterm{def:rc:market-risk-measures:hs}{simulation historique} compte 10
dépassements sur 1\,000 jours (p-valeur de Kupiec 1.00) et aucun sur les 250 derniers, vert~; la
méthode delta-normale EWMA en compte 28 (rapport de vraisemblance de Kupiec 22.0, p-valeur $3\times10^{-6}$) et 8 sur les
250 derniers, jaune avec une hausse de 0.75~; la \omterm{def:rc:market-risk-measures:fhs}{simulation historique filtrée} en compte 8 et 3, vert
(la \cref{fig:rc:market-risk-measures:backtest}). Aucun n'échoue au test d'indépendance à 5~\%.
\end{example}

\begin{omfigure}
\begin{tikzpicture}
\begin{axis}[width=11.5cm,height=5.6cm,
  xlabel={jours de bourse depuis le 2 septembre 2022},ylabel={millions de dollars},
  tick label style={font=\small},label style={font=\small},
  xmin=0,xmax=1000,ymin=-4.5,ymax=4.5,grid=major,grid style={gray!25},
  legend columns=2,legend cell align=left,legend style={font=\footnotesize,at={(0.5,-0.25)},anchor=north,draw=none,fill=none,/tikz/every even column/.append style={column sep=8pt}},
  table/col sep=comma]
\addplot[omExo!60,only marks,mark=*,mark size=0.5pt] table[x=day,y=pnl]{figdata/rates-credit-risk/21-market-risk-measures/backtest.csv};
\addplot[omDef,thick] table[x=day,y=var_hs]{figdata/rates-credit-risk/21-market-risk-measures/backtest.csv};
\addplot[omMeth,thick] table[x=day,y=var_ewma]{figdata/rates-credit-risk/21-market-risk-measures/backtest.csv};
\addplot[omThm,only marks,mark=x,mark size=2.5pt,thick] table[x=day,y=pnl]{figdata/rates-credit-risk/21-market-risk-measures/exceptions.csv};
\legend{P\&L quotidien,$-$VaR historique,$-$VaR EWMA,dépassements EWMA}
\end{axis}
\end{tikzpicture}

\omcaption{Quatre ans de P\&L quotidien du livre, face à l'opposé de la
VaR à 99~\% de la veille, issue de la \omterm{def:rc:market-risk-measures:hs}{simulation historique} et du modèle delta-normal EWMA~; les croix marquent les
dépassements du modèle EWMA. La courbe EWMA réagit vite à la volatilité mais passe trop près de zéro
lors des jours à queues épaisses. Données~: Trésor américain, BCE~; tutoriel du chapitre.}\label{fig:rc:market-risk-measures:backtest}
\end{omfigure}

\begin{dated}{2026-09}{VaR et ES dans les règles}\label{dat:rc:market-risk-measures:basel}
RiskMetrics a été lancé par J.P. Morgan en octobre 1994 (quatrième édition de son document technique,
avec Reuters, en décembre 1996), avec des facteurs de décroissance de 0.94 pour l'horizon quotidien et 0.97 pour
les horizons mensuels. Le cadre de backtest du Comité de Bâle de 1996 a fixé les feux tricolores pour la
VaR à 99~\% sur 250 jours. Sa norme 2019 sur le risque de marché a remplacé la VaR par l'\omterm{def:rc:market-risk-measures:es}{expected shortfall} à
97.5~\% sur un horizon de base de dix jours, avec des backtests au niveau du desk à 97.5~\% et 99~\%~: plus de 12
dépassements à 99~\% ou 30 à 97.5~\% en douze mois envoient un desk vers l'approche standardisée.
\end{dated}

\section{Contributions au risque}

La \omterm{def:rc:market-risk-measures:param}{VaR paramétrique} est homogène de degré un par rapport aux positions, donc l'\omterm{def:rc:the-valuation-adjustment-desk:euler}{allocation d'Euler}
(chapitre~20) la répartit en contributions qui s'additionnent~: $c_i = z_\alpha\delta_i(\Sigma\delta)_i/\sigma$.

\begin{example}[Où se situe le risque]\label{ex:rc:market-risk-measures:euler}
Sur la VaR EWMA de 1.58~million de dollars, l'obligation à dix ans contribue pour 114~\%, EURUSD pour 23~\%,
la position vendeuse à deux ans $-32\%$ et celle sur le yen $-5\%$ (\cref{fig:rc:market-risk-measures:euler})~: la
couverture à deux ans compense une partie du risque de taux de l'obligation à dix ans, et les contributions négatives signalent des couvertures.
\end{example}

\begin{omfigure}
\begin{tikzpicture}
\begin{axis}[width=11.5cm,height=5cm,ybar,bar width=20pt,
  symbolic x coords={2y yield,10y yield,EURUSD,USDJPY},xtick=data,xticklabels={{taux 2A},{taux 10A},EURUSD,USDJPY},
  ylabel={contribution (millions de dollars)},
  tick label style={font=\small},label style={font=\small},
  ymin=-0.8,ymax=2.2,grid=major,grid style={gray!25},
  nodes near coords,every node near coord/.append style={font=\footnotesize,/pgf/number format/fixed,/pgf/number format/precision=2},
  table/col sep=comma]
\addplot[fill=omDef!60,draw=omDef] table[x=factor,y=contrib]{figdata/rates-credit-risk/21-market-risk-measures/euler.csv};
\end{axis}
\end{tikzpicture}

\omcaption{Contributions d'Euler de chaque \omterm{def:rc:market-risk-measures:factor}{facteur de risque} à la VaR EWMA à 99~\% du portefeuille. Elles s'additionnent
pour donner la VaR~; les contributions négatives viennent des positions qui couvrent le reste. Données~: le
tutoriel du chapitre.}\label{fig:rc:market-risk-measures:euler}
\end{omfigure}

\section{Tutoriel~: VaR sur données publiques}\label{tut:rc:market-risk-measures:tutorial}

\begin{tutorial}
\textbf{Objectif.} Calculer la VaR et l'ES historiques, paramétriques, Monte Carlo et filtrées du portefeuille
sur les données du Trésor américain et de la BCE, les backtester sur quatre ans, et répartir la VaR. \textbf{État final~:}
les résultats des
\cref{ex:rc:market-risk-measures:methods,ex:rc:market-risk-measures:backtest,ex:rc:market-risk-measures:euler}
et les cinq graphiques.
\begin{tutsteps}
\item \textbf{Données}~: \texttt{data()} réunit les taux au pair et les taux de la BCE~; \texttt{changes()}
calcule les variations des \omterm{def:rc:market-risk-measures:factor}{facteurs de risque}.
\item \textbf{Mesures}~: \texttt{measures()} exécute chaque méthode sur les 500 derniers jours.
\item \textbf{Backtest}~: \texttt{backtest()} fait avancer les modèles sur 1\,000 jours.
\item \textbf{Graphiques}~: \texttt{fig\_rc\_var.py}.
\end{tutsteps}
\textbf{Pour aller plus loin.} Utiliser $\lambda = 0.97$~; retirer le straddle et comparer le delta et la
réévaluation complète~; raccourcir la fenêtre à 250 jours et observer la réaction de la VaR historique.
\end{tutorial}

\section{Construction~: le modèle de VaR}\label{bld:rc:market-risk-measures:varmodel}

\begin{build}
\bfield{Rôle} Les mesures de risque de marché de l'entreprise et leurs tests, utilisés par les chapitres
consacrés au stress, au capital et au moteur de risque qui suivent.
\bfield{Interface} \texttt{hs\_var\_es(pnl, alpha)}~; \texttt{ewma\_cov}, \texttt{ewma\_vol\_path},
\texttt{fhs\_scenarios}~; \texttt{parametric\_var\_es}, \texttt{euler\_var}~; \texttt{mc\_var\_es(revalue,
cov, alpha)}~; \texttt{delta\_gamma}~; \texttt{kupiec}, \texttt{christoffersen},
\texttt{traffic\_light}~; \texttt{sqrt\_time}.
\bfield{Règles} Pertes positives~; la VaR historique est la $\lceil(1-\alpha)n\rceil$-ième plus grande perte, et l'ES
la moyenne de ce même nombre~; poids EWMA normalisés~; \omterm{def:rc:market-risk-measures:traffic}{test des feux tricolores} sur 250 jours.
\bfield{Tests d'acceptation} \texttt{code/firm/varmodel/tests/}~: VaR et ES d'un échantillon connu~;
Monte Carlo contre paramétrique pour un portefeuille linéaire~; les contributions d'Euler s'additionnent~; ES gaussienne à
97.5~\% proche de la VaR à 99~\%~; \omterm{def:rc:market-risk-measures:tests}{test de Kupiec} nul au nombre attendu de dépassements~; bornes du \omterm{def:rc:market-risk-measures:traffic}{test des feux tricolores}~;
le \omterm{def:rc:market-risk-measures:tests}{test de Christoffersen} détecte les regroupements~; les scénarios filtrés intègrent la volatilité du jour.
\bfield{Pour aller plus loin} Queues de valeurs extrêmes~; filtrage GARCH~; backtests de l'ES~; calibration sous stress~;
ES par composante.
\end{build}

\begin{omsources}
\item J.P.\ Morgan/Reuters, \emph{RiskMetrics --- Technical Document}, 4e éd., 1996.
\item P.\ Artzner, F.\ Delbaen, J.-M.\ Eber et D.\ Heath, ``Coherent measures of risk'',
\emph{Mathematical Finance} 9(3), 1999, 203--228.
\item P.\ Kupiec, ``Techniques for verifying the accuracy of risk measurement models'',
\emph{Journal of Derivatives} 3(2), 1995, 73--84.
\item P.\ Christoffersen, ``Evaluating interval forecasts'', \emph{International Economic Review}
39(4), 1998.
\item Basel Committee on Banking Supervision, cadre de backtesting (1996)~; \emph{Minimum capital
requirements for market risk} (2019).
\end{omsources}

\section{Exercices}

\begin{exercise}[$\star$]\label{exo:rc:market-risk-measures:1}
Avec 500 jours de données historiques, quelle perte est la VaR à 99~\% et quelles pertes composent l'ES à 97.5~\%~?
\end{exercise}

\begin{exercise}[$\star$]\label{exo:rc:market-risk-measures:2}
Transposer la VaR EWMA d'un jour à dix jours avec la \omterm{def:rc:market-risk-measures:sqrt}{règle de la racine carrée du temps}. Dans quels cas la
règle est-elle fausse~?
\end{exercise}

\begin{exercise}[$\star$]\label{exo:rc:market-risk-measures:3}
Un modèle présente 7 dépassements en 250 jours. Donner la zone et le facteur de multiplication.
\end{exercise}

\begin{exercise}[$\star\star$]\label{exo:rc:market-risk-measures:4}
Vérifier les chiffres de l'\cref{ex:rc:market-risk-measures:subadd}.
\end{exercise}

\begin{exercise}[$\star\star$]\label{exo:rc:market-risk-measures:5}
Pourquoi le straddle vendeur rend-il la VaR en delta seul trop faible~?
\end{exercise}

\begin{exercise}[$\star\star$]\label{exo:rc:market-risk-measures:6}
Pourquoi le modèle EWMA a-t-il échoué au \omterm{def:rc:market-risk-measures:tests}{test de Kupiec} alors qu'il suivait bien la volatilité~?
\end{exercise}

\begin{exercise}[$\star\star\star$]\label{exo:rc:market-risk-measures:7}
\emph{Programmation.} Calculer la statistique de Kupiec et la valeur p pour 28 dépassements en 1\,000 jours à 1~\%.
\end{exercise}

\begin{exercise}[$\star\star\star$]\label{exo:rc:market-risk-measures:8}
\emph{Trouver la faille.} «~Notre VaR historique n'a enregistré aucun dépassement au cours des 250 derniers jours, elle est donc trop
prudente et nous devrions passer au chiffre plus bas de l'EWMA.~»
\end{exercise}

\section{Problème~: le rapport de seize heures}

\begin{problem}[{Problème du week-end --- un chiffre, et ce qu'il coûte}]\label{pb:rc:market-risk-measures:1}
Le portefeuille de l'\cref{ex:rc:market-risk-measures:book} est un desk dont le capital est fixé par son modèle EWMA
delta-normal selon le cadre de 1996~: trois fois (plus la majoration du backtesting) la VaR
à 99~\% sur dix jours, mise à l'échelle à partir d'un jour.

\medskip\noindent\textbf{Partie I --- Le chiffre.}
\begin{enumerate}
\item Donner la VaR à 99~\% à un jour et l'ES à 97.5~\% du desk, issues du modèle EWMA et de la
\omterm{def:rc:market-risk-measures:hs}{simulation historique}.
\item Quel \omterm{def:rc:market-risk-measures:factor}{facteur de risque} domine, et de combien~?
\item Donner la VaR sur dix jours.
\item Pourquoi l'ES n'est-il que légèrement supérieur à la VaR dans le modèle EWMA~?
\item Quel chiffre faudrait-il présenter au conseil d'administration, et pourquoi~?
\end{enumerate}

\medskip\noindent\textbf{Partie II --- Le backtest.}
\begin{enumerate}[resume]
\item Donner les dépassements sur les 250 derniers jours et la zone.
\item Donner le facteur de multiplication.
\item Donner le capital avec et sans la majoration.
\item Combien de dépassements la \omterm{def:rc:market-risk-measures:hs}{simulation historique} a-t-elle enregistrés la même année~?
\item La \omterm{def:rc:market-risk-measures:fhs}{simulation historique filtrée} modifierait-elle le capital~?
\end{enumerate}

\medskip\noindent\textbf{Partie III --- Le modèle.}
\begin{enumerate}[resume]
\item Pourquoi un modèle normal à volatilité rapide présente-t-il encore trop de dépassements~?
\item Qu'exigerait à la place de cette VaR la norme de 2019~?
\item Comment le desk se situerait-il par rapport à sa règle des 12 et 30 dépassements~?
\item Qu'apporte le test d'indépendance~?
\item Quel modèle faudrait-il proposer, et combien coûterait-il en capital~?
\end{enumerate}

\medskip\noindent\textbf{Partie IV --- Jugement.}
\begin{enumerate}[resume]
\item Pourquoi la VaR la plus basse n'est-elle pas le meilleur modèle~?
\item Un desk peut-il contourner une limite de VaR~? Comment~?
\item Que devrait contenir le rapport de seize heures, outre la VaR~?
\item Énoncer le \emph{résultat nommé}~: la VaR à 99~\% et l'ES à 97.5~\% du desk, ses dépassements sur 250
jours et le multiplicateur de capital qu'ils impliquent.
\item En une phrase~: que mesure la VaR et que ne mesure-t-elle pas~?
\end{enumerate}
\end{problem}

\section{Questions d'entretien}

\begin{interviewq}[$\star$ \iqroles{risque, trader}]\label{iq:rc:market-risk-measures:1}
Définissez la VaR et l'ES. Pourquoi les régulateurs sont-ils passés de l'une à l'autre~?
\end{interviewq}

\begin{interviewq}[$\star\star$ \iqroles{risque}]\label{iq:rc:market-risk-measures:2}
Comparez la \omterm{def:rc:market-risk-measures:hs}{simulation historique}, la \omterm{def:rc:market-risk-measures:param}{VaR paramétrique} et la \omterm{def:rc:market-risk-measures:mc}{VaR Monte Carlo}.
\end{interviewq}

\begin{interviewq}[$\star\star$ \iqroles{chercheur, risque}]\label{iq:rc:market-risk-measures:3}
Montrez, sur un exemple, que la VaR n'est pas sous-additive.
\end{interviewq}

\begin{interviewq}[$\star\star$ \iqroles{risque, banque}]\label{iq:rc:market-risk-measures:4}
Comment effectuez-vous le backtest d'un modèle de VaR~? Que testent les \omterm{def:rc:market-risk-measures:tests}{tests de Kupiec} et de Christoffersen~?
\end{interviewq}

\begin{interviewq}[$\star\star\star$ \iqroles{développeur}]\label{iq:rc:market-risk-measures:5}
Une VaR historique sur 100\,000 transactions avec réévaluation complète prend six heures. Comment la rendez-vous
rapide~?
\end{interviewq}

\begin{interviewq}[$\star\star\star$ \iqroles{chercheur}]\label{iq:rc:market-risk-measures:6}
Comment effectueriez-vous le backtest de l'\omterm{def:rc:market-risk-measures:es}{expected shortfall}~?
\end{interviewq}
